\section{问题重述}
\subsection{问题背景}
情感分析是具身智能、自然人机交互等领域的关键支撑技术。真实交互中，人类情绪经由语言内容、语音韵律与面部运动等多通道信号协同表达：词义决定评价的正负方向，语调、语速与能量承载强调、迟疑与情绪唤醒，面部关键点的运动则补充词面之外的态度信息。三者的物理载体并不相同，分别表现为离散词元序列、等间隔采样的声学帧与带时间戳的视频帧，在采样频率、特征维度与时间尺度上差异显著。多模态联合建模因此成为突破单模态识别局限、提升复杂场景下情绪理解准确性与稳定性的核心路径\cite{ref:mm}。

然而，多模态情感预测在真实场景中面临三重挑战。其一，特征提取与时序对齐的质量直接决定模型性能上限：自动转写可能识别失败，音视频轨道的起止时刻与有效时长并不一致，部分视频帧中人脸检出失败，若按序号直接拼接三路特征，同一位置上的语义、语音与视觉信息可能对应不同时刻，若直接剔除不完整样本，又会损失本就有限的建模数据。其二，画面遮挡、环境噪声与转写失败常导致模态局部时段不可用，常规固定权重融合模型的性能随之急剧下降。其三，多数模型仅输出预测结果，决策依据难以追溯与复核，制约了模型的可信应用。因此，多模态情感预测需要在三个层次上开展建模：先建立从原始素材到标准化时序特征的对应关系，再在证据不完整时稳定输出情感极性与情感强度，最后量化说明判断所依据的模态与局部片段。

\subsection{问题描述}
本文以赛题提供的CMU-MOSEI系列数据\cite{ref:mosei}为唯一数据来源，依次完成三项建模任务。

\textbf{问题一：}以附件1的原始视频为研究对象，建立文本、语音、视觉三模态情感特征提取与时序对齐的数学模型。模型须明确三类情感特征的定义与提取方法，给出跨模态时间对应规则，将采样频率与时间粒度各异的三路信号统一至同一时间基准，并规定各处理环节对原始样本的操作方式。自生成的特征文件需满足三项规范：样本编号、模态文件与输出特征一一对应；序列位置、有效长度、填充规则及其与原始素材的对应记录完整可核验；特征提取工具、关键参数、处理日志与复现实验说明齐备。最终以表格形式汇总附件1全部100条样本的特征提取结果，并以典型样本呈现文本片段、语音时段、视频帧段与三类特征的对应关系。

\textbf{问题二：}针对模态局部缺失的复杂场景，建立鲁棒性多模态情感预测模型。模态局部缺失指一条或多条模态在部分连续时段内信息不可用，而非整条模态缺失。模型须在模态信息不完整的条件下稳定输出情感极性与情感强度预测，并分析缺失模态类型、缺失位置、缺失时长及缺失比例对预测性能的影响规律。模型建立后须对附件3中的测试样本完成推理预测。

\textbf{问题三：}针对三模态信息完整而决策依据难以追溯的场景，建立可解释性多模态情感预测模型。模型除输出情感极性与情感强度外，还须以可量化、可复核的方式给出判断依据：量化各模态在当前样本中的作用差异，明确主要参考模态，定位各模态中与判断密切相关的局部片段。关键证据应可对应至原始文本片段、语音时段或视觉关键帧，以支持回溯核验。模型须对附件4中的测试样本完成预测与解释。

\subsection{总体解题思路}
三项任务构成从时间组织到鲁棒预测、再到决策解释的递进链条。问题一处理原始素材与时间基准：三模态采样粒度不同，需先以词为公共时间锚点建立可核验的对应关系，输出内容位置、观测掩码与词级时间区间。问题二处理证据不完整：在附件2的官方对齐特征上区分填充、观测与局部缺失三种状态，使预测在连续缺口存在时仍然可用。问题三处理决策可追溯：在附件2上独立训练预测器，通过固定参考类别的模态联盟分解与输入片段删除，给出主要参考模态、模态作用程度与可回映到原始素材的关键证据。

三问的数据输入、模型节点与输出关系如图\ref{prob:fig:framework}所示。

\begin{figure}[htbp]
\centering
\resizebox{\linewidth}{!}{%
\begin{tikzpicture}[
  box/.style={draw=black!45,fill=black!3,rounded corners=2pt,
    text width=4.55cm,minimum height=1.12cm,align=center,font=\small},
  arrow/.style={-{Latex[length=2mm]},thick,draw=black!65}]
\node[font=\small\bfseries] at (0,0.55) {问题一：特征与时间组织};
\node[font=\small\bfseries] at (5.2,0.55) {问题二：局部缺失预测};
\node[font=\small\bfseries] at (10.4,0.55) {问题三：贡献与证据解释};
\node[box,fill=blue!6] (a1) at (0,-0.45) {附件1：100条原始视频\\官方转写与音视频信号};
\node[box,fill=green!6] (b1) at (5.2,-0.45) {附件2：训练与验证\\附件3：专项预测};
\node[box,fill=orange!7] (c1) at (10.4,-0.45) {附件2：训练与验证\\附件4：预测与证据回映};
\node[box] (a2) at (0,-2.05) {文本、语音与视觉特征\\CTC原词区间与子词归属};
\node[box] (b2) at (5.2,-2.05) {三态掩码与观测节点补全\\可靠度门控与双向时序编码};
\node[box] (c2) at (10.4,-2.05) {模态子集训练与双任务预测\\固定参考类别的联盟分解};
\node[box] (a3) at (0,-3.65) {时间交叠池化\\覆盖率与边界扰动分析};
\node[box] (b3) at (5.2,-3.65) {标签、蒸馏与重建联合目标\\缺失因素分析与组件消融};
\node[box] (c3) at (10.4,-3.65) {局部删除与忠实性检验\\种子及解释基线敏感性};
\node[box,fill=blue!6] (a4) at (0,-5.25) {100条自主特征与时间记录\\三模态维度：768/25/16};
\node[box,fill=green!6] (b4) at (5.2,-5.25) {附件3全部30条预测\\情感极性与连续强度};
\node[box,fill=orange!7] (c4) at (10.4,-5.25) {附件4全部20条预测与解释\\模态贡献与原始关键片段};
\foreach \prefix in {a,b,c} {
  \draw[arrow] (\prefix1)--(\prefix2);
  \draw[arrow] (\prefix2)--(\prefix3);
  \draw[arrow] (\prefix3)--(\prefix4);
}
\end{tikzpicture}
}
\caption{本文总体研究思路}
\label{prob:fig:framework}
\end{figure}

三问共享时间组织与掩码表示方法，各自使用相应附件的数据。问题一生成768/25/16维自主特征；问题二与问题三使用768/74/35维官方特征，分别学习预测参数。问题三沿用词级对齐方法，为附件4建立独立的时间记录，将特征层证据回映至原始素材。
